%% kapa: exam=1 round=20 year=2009 period=1 form=A subject=civil_law subject_ko=민법 seq=37 number=37
%% source: Q-Net 감정평가사 기출문제 (cst00309) artlSeq=5024178 file=extracted/1008/5024178/2009년 제20회 감정평가사 1차시험 기출문제/1교시.hwp
%% conversion: kordoc 4.24.51 -> tools/gen1.py (auto); review=proofread-sonnet; flags=none
\begin{question}{37}
소멸시효의 기산점에 관한 설명 중 옳은 것($\bigcirc$)과 틀린 것($\times$)을 바르게 표시한 것은? (다툼이 있으면 판례에 의함)
\begin{framedtext}
ㄱ. 권리가 확정기한부라면 소멸시효의 기산점은 그 기한이 도래한 때이고, 불확정기한부인 경우에는 기한이 객관적으로 도래한 때이다. \par
ㄴ. 부작위를 목적으로 하는 채권의 소멸시효는 위반행위를 한 때부터 진행한다. \par
ㄷ. 부당이득반환청구권은 그 성립과 동시에 행사할 수 있으므로, 원칙적으로 그때부터 소멸시효가 진행한다. \par
ㄹ. 동시이행의 항변권이 붙어 있는 채권은 상대방이 이행의 제공을 한 때로부터 소멸시효가 진행한다. \par
\end{framedtext}
\begin{choices}
  \item ㄱ($\bigcirc$) ㄴ($\bigcirc$) ㄷ($\bigcirc$) ㄹ($\times$)
  \item ㄱ($\bigcirc$) ㄴ($\bigcirc$) ㄷ($\times$) ㄹ($\bigcirc$)
  \item ㄱ($\times$) ㄴ($\bigcirc$) ㄷ($\times$) ㄹ($\times$)
  \item ㄱ($\bigcirc$) ㄴ($\times$) ㄷ($\bigcirc$) ㄹ($\bigcirc$)
  \item ㄱ($\times$) ㄴ($\times$) ㄷ($\bigcirc$) ㄹ($\bigcirc$)
\end{choices}
\end{question}
\answer{1}
